\chapter{Authentication And Trust Boundaries}

\begin{longtable}{@{}L{0.22\textwidth}L{0.24\textwidth}L{0.22\textwidth}L{0.22\textwidth}@{}}
\caption{Actor and authentication matrix}\\
\toprule
\textbf{Actor} & \textbf{API surface} & \textbf{Auth model} & \textbf{Scope source} \\
\midrule
\endfirsthead
\toprule
\textbf{Actor} & \textbf{API surface} & \textbf{Auth model} & \textbf{Scope source} \\
\midrule
\endhead
Merchant backend & Merchant API & HMAC headers: merchant id, access key, timestamp, signature. & Authenticated merchant credential. \\
Provider simulator & Provider callback API & Environment-local trust for MVP simulator routes. & Callback payload references payment or refund records. \\
Admin/Ops user & Ops API & Internal HttpOnly session cookie. & Internal user role \apiid{ADMIN} or \apiid{OPS}. \\
Admin user & Portal-user provisioning API & Internal HttpOnly session cookie plus \apiid{ADMIN} role. & Internal user role. \\
Merchant portal user & Merchant Portal API & Merchant HttpOnly session cookie. & \texttt{merchant\_users.merchant\_db\_id}. \\
\bottomrule
\end{longtable}

Important trust-boundary rules:

\begin{itemize}
  \item Merchant Portal routes never accept a client-supplied
  \texttt{merchant\_id} for data scope after login.
  \item Internal write routes still accept a nested actor object so the
  user-supplied reason can become part of the audit trail, but actor identity
  is normalized from the authenticated internal session.
  \item Provider callback routes are simulator-trusted in the MVP and should be
  treated as internal or controlled-environment routes rather than public,
  internet-grade provider integrations.
\end{itemize}
